% page 367 of the facsimile (image p-366 of the scan)
% block 162.6 x 242.0 mm ; left margin 39.4 mm
\pgc{17.780mm}{0.000mm}{
\l{\cel{53}359}
\l{}
\l{\sou{martiro}.- I. La persono qua tormentesis od ocidesis pro ne}
\l{\cel{3}renegar sua profeso di fido a Iesu, en la historio di krist-}
\l{\cel{3}anismo. - II. La persono qua supliciesis sen agir por jus-}
\l{\cel{3}tifikar lo. - III. La persono qua sufris por permanar fidela}
\l{\cel{3}a sua idealo. - DEFIS.}
\l{}
\l{\sou{martro}. (\sou{zool.}) Mikra mamifero di qua la korpo esas longa e la}
\l{\cel{3}felo tre prizata. - L.\cel{1}\sou{mustela martes}. DEFIS.}
\l{}
\l{\sou{marvelo}. To quo (o : ta qua) astonas pro lua beleso, lua}
\l{\cel{3}grandeso. - EFIS.}
\l{}
\l{\sou{maso}.\cel{2}Asemblajo grandega de kozi, o parti de kozi, qui formacas}
\l{\cel{3}quaza ensemblo.\cel{2}- DEFIRS.}
\l{}
\l{\sou{masajar}. (\sou{trans.)}\cel{2}Palpar, klemar, petrisar per manui (la parti}
\l{\cel{3}diversa di la homo-korpo) por acelerar la sango-cirkulo, por}
\l{\cel{3}igar plu flexebla la membri. - DEFIRS.}
\l{}
\l{\sou{masakrar.}\cel{1}(\sou{trans}.) Ocidar per frapar obstinoze, e pelmele. -}
\l{\cel{3}DEFIS.}
\l{}
\l{\sou{masho}. I. Singla ek la mikra slingi de filo ek lino, kotono,}
\l{\cel{3}silko, kordo, di qui la interplekto formacas texuro laxa}
\l{\cel{3}(reto, trikoturo, e c.).- II. Spaco vakua, di qua ta slingo}
\l{\cel{3}formacas la konturo. - DE.}
\l{}
\l{\sou{mashino}. Asembluro de peci, kombinita tale ke li produktas ula}
\l{\cel{3}efekti. DEFIRS.}
\l{}
\l{\sou{mashinaca}r.(\sou{trans.)}\cel{2}Kombinar artificoze, ruzoze, ula moyeni por}
\l{\cel{3}atingar skopo quan onu ne audacas agnoskar publike. - DEFIS.}
\l{}
\l{\sou{masiva.}\cel{2}Qua esas maso kompakta. - DEFIRS.}
\l{}
\l{\sou{maskar}. (\sou{trans}.)Celar sub ta peco de veluro o de satino, di qua}
\l{\cel{3}la formo quaze muldas la vizajo, kun trui qui korespondas}
\l{\cel{3}a la du okuli ed a la boko - e quan olim +weris la mulieri}
\l{\cel{3}por prezervar sua karnaciono - e quan eli +weras ankore nun}
\l{\cel{3}en la bali de maskilieri. - DEFIRS.}
\l{}
\l{\sou{maskarono}. (\sou{arkitekt.}) Figuro di fantazio, alta-reliefa, o}
\l{\cel{3}basreliefa, per qua onu ornas entablamenti, klozo-petri. -}
\l{\cel{3}DEFRS.}
\l{}
\l{maskulo.\cel{2}I. Individuo ek la sexuo qua havas la fakultato}
\l{\cel{3}fekundigar la altra sexuo. - II.\cel{1}\sou{maskula}\cel{1}(gram.): Qua aparte-}
\l{\cel{3}nas a la genro quan onu uzas por indikar la enti maskula, o}
\l{\cel{3}ti qui (arbitriale) esas asimilita ad ilci. - DEFIS.}
\l{}
\l{masonar. (trans.)\cel{3}(en\cel{1}\sou{la konstrukto di edifico, di domo,}\cel{1}di}
\l{\cel{3}\sou{ponto, l}\cel{1}o.), Facar ta parto di la laboruro qua konsistas ye}
\l{\cel{3}la asemblo di la quadri, di la kruda petri, di la briki, e c.}
\l{\cel{3}e la junturi, la induturi ek mortero, gipao, e c. - EF.}
}
